\section{The projective monomial source group}

\subsection{Finite reduction}

Suppose a monomial transformation
\(x_i\mapsto a_ix_{\sigma(i)}\) sends \(C\) to \(\alpha C\).
Coefficient comparison gives \(a_i^3=\alpha\) for all \(i\).  After a
common projective rescaling, it therefore has the form
\[
 x_i\longmapsto \rho^{e_i}x_{\sigma(i)},\qquad
 e_i\in\F_3,\quad e_0=0.                            \tag{3.1}
\]
The support of \(Q\) is an eight-cycle, so \(\sigma\) lies in its
order-\(16\) dihedral automorphism group.  Write
\[
 r_k(i)=i+k,\qquad s_k(i)=k-i\pmod8.
\]

Let \(c(E)\in\F_3\) be zero on the seven ordinary edges and one on the
closing edge \(\{7,0\}\).  Preservation of \(Q\) up to
\(\rho^{q_g}\) is equivalent to
\begin{equation}
\label{eq:phase-system}
 c(\{i,i+1\})+e_i+e_{i+1}
 =q_g+c(\{\sigma(i),\sigma(i+1)\})
\end{equation}
for all cyclic indices \(i\).  This is an affine linear system over
\(\F_3\).

\begin{proposition}
\label{prop:enum}
The only support permutations admitting solutions to
\eqref{eq:phase-system} are
\[
 r_0,r_2,r_4,r_6,s_1,s_3,s_5,s_7.
\]
Each admits exactly three normalized phase vectors.  Hence
\(|\Gm|=24\).
\end{proposition}

\begin{proof}
Row reduction of \eqref{eq:phase-system} for the sixteen possible
support permutations yields the complete table in
Appendix~\ref{app:tables}.  Substitution verifies every surviving row,
and the exhausted support list proves completeness.
\end{proof}

\subsection{Presentation}

Choose
\[
\begin{aligned}
r:&\quad \sigma=(6,7,0,1,2,3,4,5),&
e&=(0,1,1,0,1,0,1,0),\\
s:&\quad \sigma=(7,6,5,4,3,2,1,0),&
e&=(0,1,0,1,0,1,0,1).
\end{aligned}                                      \tag{3.2}
\]
Direct substitution gives
\[
 C(rx)=C(sx)=C(x),\qquad Q(rx)=Q(sx)=\rho Q(x).
\]
Exact phase-permutation multiplication gives
\[
 r^{12}=s^2=1,\qquad srs=r^{-1}.                   \tag{3.3}
\]
The twelve \(r^k\) and twelve \(sr^k\) are distinct.  By
Proposition~\ref{prop:enum}, they exhaust \(\Gm\).  This proves
\[
 \Gm\cong\Dih(C_{12})=C_{12}\rtimes C_2.
\]
We always specify order \(24\), since dihedral-group notation is not
uniform across the literature.

\begin{remark}
The proof uses the particular embedding of the two equations in
\(\PP^7\).  It does not rule out nonmonomial automorphisms of \(X\).
\end{remark}
